\documentclass{situationist-preprint}

\title{Adumbrations Do Not Converge}
\msid{PHE-03-a}
\msfamily{PHE-03}
\msdiscipline{Phenomenology}
\mstype{Principal manuscript}
\mskeywords{adumbration; Abschattung; perception; horizon; fulfillment; Husserl; Merleau-Ponty; Sartre; Noë}

\begin{document}
\maketitle

\begin{abstract}
Phenomenology began with the demand to return to the things themselves. Can the demand be met? We argue that it can, but never all at once. Every spatial thing is given in adumbrations: profiles that present the thing itself, from a side, under conditions. We distinguish three questions: whether an adumbration gives the thing itself, whether every determination of the thing can be given in some adumbration, and whether the series of adumbrations converges on a complete givenness, in the sense that from some stage on every determination is held in fullness. From four principles (self-givenness, horizonal accessibility, the openness of the inner horizon, and the finitude of what one phase of perception can hold in fullness) we prove that the first two questions have a positive answer and the third a negative one: no stage of any series gives the thing completely, and no series converges, since infinitely many determinations would have to be held at once by a bounded capacity. We reconstruct the history of the problem from Husserl's \emph{Thing and Space} to recent sensorimotor and fulfillment theories, examine the cube, the melody and the other person as cases, and answer five objections, including Husserl's own appeal to complete givenness as an Idea in the Kantian sense.
\end{abstract}

\section{The question}

In the introduction to the second volume of the \emph{Logical Investigations}, Husserl declared that the investigation of meaning must go back from words to the things themselves, to the intuitions in which meanings are fulfilled \citep[vol.~2, Introduction, §2]{Husserl1901}. The phrase became the motto of the movement. Heidegger glossed it as the maxim of phenomenology, opposed to free-floating constructions and to concepts taken over without demonstration \citep[§7]{Heidegger1927}. Merleau-Ponty described the return to the things themselves as a return to the world prior to knowledge, of which knowledge always speaks \citep[Preface]{MerleauPonty1945}.

The motto raises a question that its proponents did not always separate from its methodological use. Is a return to the things themselves possible at all? The question is pressing because of a fact that Husserl was the first to analyze in detail: a spatial thing is never given all at once. It shows itself from one side, under one illumination, at one distance, and its other sides are co-intended without being given. Husserl called the side-wise presentations of the thing its adumbrations (\emph{Abschattungen}) \citep[§§41--44]{Husserl1913}. If a thing is only ever given in adumbrations, then either the series of adumbrations, pursued far enough, delivers the thing completely, or we never reach the thing itself and have only its profiles.

Three answers are available. On the \emph{completion} view, the thing is the limit of the series of its adumbrations, and the series converges on it: enough perceptual exploration delivers the thing in full, at least in principle. On the \emph{withdrawal} view, the thing itself stands behind its adumbrations, which are appearances of something that never appears; we return only to appearances. On the view we defend, we can return to the things themselves, since each adumbration gives the thing itself and not a representative of it, and every determination of the thing can be given in some adumbration; but we never return all at once, since the series does not converge on a complete givenness. The thesis is close to Husserl's own considered view, but it rejects one reading of his doctrine that complete givenness is an Idea in the Kantian sense.

Section~\ref{s:history} reconstructs the history of the problem. Section~\ref{s:framework} gives the framework, and Section~\ref{s:argument} the argument. Section~\ref{s:cases} examines further cases, Section~\ref{s:recent} recent literature, and Section~\ref{s:objections} objections.

\section{History of the problem}\label{s:history}

\subsection{Husserl: from Thing and Space to Ideas I}

In the \emph{Logical Investigations} Husserl distinguished adequate from inadequate perception. In adequate perception the object is given as it is in itself, with nothing intended that is not also given; in inadequate perception the object is intended as more than what is given, and the surplus is an intention awaiting fulfillment \citep[vol.~2, Sixth Investigation, §§37--39]{Husserl1901}. Spatial things can only be perceived inadequately. The point was developed in the lectures of 1907 published as \emph{Thing and Space}, where Husserl analyzed the constitution of the spatial thing through the correlation of visual and tactile adumbrations with kinaesthetic sensations: the course of profiles is motivated by movements of the eyes, head and body, so that if I move thus, the thing will appear thus \citep{Husserl1973}. The thing is constituted as the identical unity of an open manifold of such motivated courses.

In the \emph{Ideas} the analysis becomes a principle of essence. A spatial thing can be given only in adumbrations, and this is not a limitation of human perception but belongs to the essence of thing-perception: even God, Husserl says, would have to perceive things through appearances, since a thing given otherwise would not be a spatial thing \citep[§§42--44]{Husserl1913}. In contrast, a lived experience does not adumbrate itself: reflection on it grasps it as it is, though only in a phase of its flow. The thing is therefore given with an essential one-sidedness, and its perception always contains empty intentions pointing to further sides.

Later in the same book Husserl asks whether the thing could, in principle, be given completely. His answer is that complete givenness of a thing is prescribed as an Idea in the Kantian sense: not attainable in any finite course of experience, but determinate as the rule of the infinite continuum of appearances in which the thing would be given \citep[§§142--143]{Husserl1913}. The Idea is itself given with evidence, as the Idea of an infinity that no finite process exhausts. This doctrine has been read in two ways. On one reading it endorses the completion view in a regulative form: the series converges, though the limit is never reached in finite time. On the other it denies convergence and treats the Idea only as the rule that organizes the series. We argue for the second reading in Section~\ref{s:objections}.

\subsection{Gurwitsch and the noematic system}

\citet{Gurwitsch1964} combined Husserl's theory with Gestalt psychology. The perceptual noema, the thing as perceived from a side, is an organized whole in which the given side and the co-intended sides are functionally interdependent. The thing itself is the system of all its noemata, the group of all the ways it can appear, held together by the identity that each noema refers to. On this account, to perceive the thing itself is to perceive one member of the system as a member of it. Gurwitsch's view comes close to the completion view in its identification of the thing with the system of appearances, but he did not hold that the system could be given at once. The system is given only as implied in each of its members.

\subsection{Sartre: the infinite in the finite}

Sartre's introduction to \emph{Being and Nothingness} begins with the claim that modern thought has reduced the existent to the series of appearances that manifest it \citep{Sartre1943}. This reduction, Sartre argues, does not abolish the distinction between appearance and being; it transforms it into the distinction between the finite and the infinite. Each appearance is finite, but the object is the principle of the series of its appearances, and the series is infinite. The object is therefore transcendent to each appearance and even to any finite collection of them, and it is, in Sartre's phrase, the infinite in the finite. Sartre further argued that the being of the phenomenon is not reducible to the phenomenon of being: the series of appearances requires a transphenomenal being that grounds it. We do not follow Sartre in this last step, but his formulation of the problem is the one we adopt: the question is whether the infinite series of appearances can be given, and the answer is that it cannot be given at once.

\subsection{Merleau-Ponty: the house seen from everywhere}

Merleau-Ponty began from the same analysis and gave it a bodily turn. To see an object is to inhabit it and to grasp from there all things in their visible aspects; the house itself is not the house seen from nowhere but the house seen from all perspectives at once, and that totality is not a further perspective that anyone occupies \citep[Part One, Introduction]{MerleauPonty1945}. The synthesis of perspectives is not an intellectual act but the work of the body schema, which holds the other sides as possibilities of movement. Merleau-Ponty insisted that the thing remains inexhaustible: it is given as more than any perception of it, and this excess belongs to its reality. The thing \enquote{resists} us; a thing that could be given completely would be an object of thought rather than of perception. His later work radicalized this point: perceptual faith is the conviction that the thing is there, beyond each of its profiles, and this conviction is never cashed out in a complete perception \citep{MerleauPonty1964}.

\subsection{Withdrawal views}

The withdrawal view has two sources. The first is the Kantian distinction between appearances and things in themselves, on which what is given is always an appearance and the thing in itself is not given at all. The second is a reading of Husserl's own analyses in which adumbrations are sensuous contents, the hyletic data of the \emph{Ideas}, which represent the thing as signs represent their objects. Husserl explicitly rejected the second reading: the thing is not given through an image or a sign but in person (\emph{leibhaftig}), and to treat perception as the apprehension of signs is to confuse perception with imagination or signification \citep[§§43, 52]{Husserl1913}. The first source we set aside, since the Kantian thing in itself is not a spatial thing given in perception and so not the thing to which the phenomenological motto bids us return.

\section{Framework}\label{s:framework}

\begin{definition}\label{d:adumbration}
An \emph{adumbration} of a spatial thing is a presentation of the thing in which some of its determinations are given intuitively, in a determinate orientation and under determinate conditions, and the remaining determinations are co-intended as belonging to the same thing without being given.
\end{definition}

\begin{definition}\label{d:determinable}
A \emph{determination} of a thing is a respect in which it can be given: the color of its rear face, the texture of its surface at a given scale, its weight in the hand, its sound when struck, its appearance under raking light. A determination is given \emph{in fullness} at a phase of perception if it is either presented in the current adumbration or still held in retention with the intuitive content it had when presented.
\end{definition}

\begin{definition}\label{d:series}
A \emph{perceptual series} is a sequence of phases $1, 2, 3, \dots$ of an exploration of one thing, each with its current adumbration. A series gives the thing \emph{completely at phase $n$} if every determination of the thing is given in fullness at $n$. It \emph{converges on a complete givenness} if every determination is, from some phase on, given in fullness at all later phases.
\end{definition}

The notion of convergence in Definition~\ref{d:series} is the weakest that the completion view requires. It does not demand that any phase give the thing completely; it demands only that each determination, once secured, be kept, so that the givenness of the thing increases without loss and in the limit includes every determination. If even this notion fails, then no stronger notion can hold.

\begin{principle}[Self-givenness]\label{p:self}
In each adumbration, the thing itself is given, not a representative of it.
\end{principle}

\begin{principle}[Horizonal accessibility]\label{p:access}
Every determination of a thing can be given in some adumbration of it, reachable by a course of exploration from any other.
\end{principle}

\begin{principle}[Open inner horizon]\label{p:open}
A spatial thing has infinitely many determinations.
\end{principle}

\begin{principle}[Finitude of the phase]\label{p:finite}
There is a bound on the number of determinations that a single phase of perception can give in fullness.
\end{principle}

Principle~\ref{p:self} is Husserl's rejection of image and sign theories of perception. Principle~\ref{p:access} is the content of the horizon: what is co-intended is co-intended as perceivable, and the kinaesthetic motivations analyzed in \emph{Thing and Space} specify the course of movement that would bring it to givenness. Principle~\ref{p:open} formulates the openness of the inner horizon. Every determination given at one scale is determinable further at another; every color is a color under an illumination, and there are indefinitely many illuminations; every surface is determinable at finer grain; and the thing's causal properties, its ways of behaving under pressure, heat and impact, form an open manifold of determinations that Husserl treated as belonging to the materiality of the thing \citep{Husserl1952}. Principle~\ref{p:finite} combines two facts. An adumbration presents the thing from one orientation, and what is presented from one orientation excludes what is presented from opposed orientations: the front and the back of an opaque thing are not presented together. And retention holds past adumbrations in fullness only within a bounded retentional horizon, beyond which they are held emptily or not at all; the living present has a finite thickness fixed by that horizon (PHE-05). A bounded presentation together with a bounded retention gives a bound on what is held in fullness at a phase.

\section{The argument}\label{s:argument}

\begin{lemma}\label{l:return}
In every phase of every perceptual series, the thing itself is given.
\end{lemma}

\begin{proof}
Each phase has a current adumbration, and by Principle~\ref{p:self} the thing itself is given in it.
\end{proof}

\begin{lemma}\label{l:each}
For every determination of a thing, there is a perceptual series in which it is given in fullness at some phase.
\end{lemma}

\begin{proof}
By Principle~\ref{p:access}, the determination can be given in some adumbration reachable by exploration from the adumbration with which a series begins. The series that follows that course of exploration contains a phase whose current adumbration presents it.
\end{proof}

\begin{lemma}\label{l:never}
No phase of any perceptual series gives the thing completely.
\end{lemma}

\begin{proof}
By Principle~\ref{p:finite}, at most $k$ determinations are given in fullness at any phase, for some bound $k$. By Principle~\ref{p:open}, the thing has more than $k$ determinations. So at every phase some determination is not given in fullness.
\end{proof}

\begin{theorem}\label{t:main}
No perceptual series converges on a complete givenness of the thing.
\end{theorem}

\begin{proof}
Suppose a series converges. Let $k$ be the bound of Principle~\ref{p:finite}, and by Principle~\ref{p:open} choose $k+1$ distinct determinations $d_1, \dots, d_{k+1}$. By convergence, for each $d_i$ there is a phase $N_i$ from which on $d_i$ is given in fullness. Let $N$ be the largest of $N_1, \dots, N_{k+1}$. At phase $N$, all $k+1$ determinations are given in fullness, contradicting the bound.
\end{proof}

\begin{corollary}\label{c:main}
We can return to the things themselves, but never all at once.
\end{corollary}

\begin{proof}
By Lemma~\ref{l:return}, every perception of a thing is a return to the thing itself, and by Lemma~\ref{l:each}, no determination of the thing is beyond the reach of such a return. By Lemma~\ref{l:never} and Theorem~\ref{t:main}, there is no phase in which all determinations are given, and no limit in which they are all held.
\end{proof}

Theorem~\ref{t:main} is stronger than Lemma~\ref{l:never}, and the difference matters. A defender of the completion view can concede that no phase gives the thing completely and still hold that the series converges, so that the givenness of the thing increases towards completeness as a sequence of partial sums increases towards its limit. The theorem excludes this. The reason is not that the series is too short; it is that what the series gives is not accumulated. A determination given at phase $n$ is held in fullness only while it remains within the retentional horizon or within the current adumbration, and beyond that it is held only emptily, as already known. The series is not a sum but a moving window.

Two further remarks clarify the structure of the result. First, Lemma~\ref{l:each} and Theorem~\ref{t:main} are compatible: each determination can be reached by some series, and indeed one series can reach each determination at some phase, if it enumerates them; what no series does is hold them. Second, the bound in Principle~\ref{p:finite} need not be fixed across subjects or circumstances. A trained observer may hold more in fullness than an untrained one. The argument requires only that the bound be finite for each series.

\begin{remark}\label{r:empty}
What is no longer given in fullness is not thereby lost. It is held emptily, as known, and can be called up again by a suitable movement. This is why the thing is given as familiar, and why its earlier sides contribute to the sense with which it is given at each phase. But empty holding is not givenness in fullness, and complete givenness, in the sense in which the completion view uses the term, requires the latter.
\end{remark}

\subsection{The intersubjective extension}

Husserl held that the full objectivity of the thing is constituted only intersubjectively: the thing is the same for all, and what I do not see from here, another sees from there \citep{Husserl1952}. This suggests a way around Theorem~\ref{t:main}. Perhaps no single series converges, but the community of perceivers, each exploring the thing from a different side, jointly gives it completely. We show that the extension does not change the result.

\begin{definition}\label{d:joint}
A \emph{joint series} of $m$ perceivers is a sequence of phases in each of which each of the $m$ perceivers has an adumbration of the same thing. A determination is \emph{jointly given in fullness} at a phase if it is given in fullness to at least one of them.
\end{definition}

\begin{proposition}\label{p:joint}
For every finite number $m$ of perceivers, no joint series converges on a complete givenness of the thing.
\end{proposition}

\begin{proof}
By Principle~\ref{p:finite}, each perceiver holds at most $k_i$ determinations in fullness at any phase, so the $m$ perceivers jointly hold at most $K = k_1 + \dots + k_m$. The proof of Theorem~\ref{t:main} then applies with $K$ in place of $k$: choose $K+1$ determinations, take the latest phase from which each is held, and derive a contradiction.
\end{proof}

Proposition~\ref{p:joint} does not deny that intersubjectivity enriches the givenness of the thing. Another's report, or her perception given to me non-primordially through her expression, lets me co-intend the side she sees with a determinacy that my own empty intention lacked. But joint givenness is bounded for every finite community at every phase, and the community of perceivers at any time is finite. An infinite community, to which the thing would be given completely, is a further Idea in the Kantian sense, which regulates the constitution of objectivity without being realized in it. The same holds for instruments: a microscope extends the range of accessible determinations, so that Principle~\ref{p:access} applies to more of them, but it gives one adumbration at a time and is subject to the bound.

The proposition also explains a familiar feature of scientific practice, that agreement about a thing is reached on determinations selected in advance (its mass, its length, its chemical composition), not on the thing in all its determinations. The selection makes a finite set of determinations the target of joint exploration, and for such a set convergence is possible, as in the case of the short melody below. The thing itself, with its open horizon, is not such a target.

\section{Cases}\label{s:cases}

\subsection{The cube}

A die on the table appears with three faces visible at most. Its six faces, one might say, are finitely many and can be given in a single rotation; so Principle~\ref{p:open} fails for the cube at the level that matters, and the cube is given completely after one turn of the hand. The objection confuses the faces with the determinations. Each face is determinable at indefinitely many scales: its pips, their depth, the grain of the material, the micro-chips at its edges; its color is determinable under indefinitely many illuminations; and its geometric properties, such as the equality of its edges, are given only to the accuracy of each perception. Even restricted to the six faces, moreover, the faces are not all given together: rotation brings three into view by removing three others, and those held in retention are held in fullness only within the retentional horizon. The cube is the simplest case of non-convergence, not a counterexample to it. The companion manuscript PHE-03-b analyzes the back of the cube in detail.

\subsection{The melody}

A melody is not a spatial thing, but it is given in temporal profiles: each phase presents a note in primal impression, with the preceding notes in retention and the expected notes in protention \citep{Husserl1966}. A short melody, unlike a spatial thing, can be held whole within the living present, if its extent falls within the retentional horizon. In that case it is given completely, and the series of its phases converges on that givenness. A longer melody, or a symphony, exceeds the horizon, and its earlier movements are held emptily while the later ones sound. Temporal objects thus illustrate both sides of the theorem: finiteness of the object together with a sufficient bound permits completeness; openness of the object or insufficiency of the bound excludes it. Spatial things are always in the second situation, because their determinations are infinitely many.

\subsection{The other person}

Another person is given in a bodily appearance, and her experiences are given through it, not primordially as they are given to her but non-primordially, as hers. If PHE-12 is right, this non-primordial givenness suffices for knowledge of other minds. It does not converge on a complete givenness either: the other's experiences are given in profiles of expression, each partial, and her experiential life has the same open horizon as a thing, with the further feature that it is given primordially only to her. The case therefore satisfies all four principles, with self-givenness (Principle~\ref{p:self}) taking the form that it is the other herself who is given, not an inference to her. The non-convergence of the series of expressions is one reason that understanding another person is an open task.

\subsection{The lived body}

The perceiver's own body is a limiting case. It is given, through proprioception and touch, in ways that are not adumbrations of the visual kind, but as an object of vision and touch it is given in a series that is restricted in a further way: certain of its determinations, such as the back of the head, are not accessible from the zero-point of orientation without instruments (PHE-08). Principle~\ref{p:access} holds for the lived body only in a weakened form. The result for the lived body is therefore stronger than Theorem~\ref{t:main}: not only does the series of its adumbrations fail to converge, but some of its determinations cannot be given to the subject as an object at all.

\section{Recent literature}\label{s:recent}

The theory of adumbrations has been revived in three bodies of recent work.

The first is the sensorimotor approach of \citet{Noe2004}. Noë asks how the whole of a tomato can be perceptually present when only its facing side is seen, and answers that it is present as accessible: the perceiver has practical knowledge of how her sensory stimulation would change with movement, and the hidden sides are present in virtue of this knowledge. Noë's account is a version of Principle~\ref{p:access}, and his \enquote{presence in absence} is what we have called empty co-intending. It sometimes suggests the completion view, since sensorimotor knowledge could seem to make the whole thing present now. But Noë himself distinguishes presence as access from presence as content: the hidden side is present as available, not as seen. Our theorem then applies without change: availability is not givenness in fullness, and the determinations held as available are not given in fullness all at once.

The second is the phenomenology of perceptual anticipation. \citet{Madary2017} argues that visual perception is a continuous cycle of anticipation and fulfillment, in which anticipations are indeterminate and become more determinate as they are fulfilled. On his account, the determinacy of anticipations increases with exploration but does not reach full determinacy; each fulfillment opens further anticipations. This is a detailed defence of Principle~\ref{p:open} in the vocabulary of predictive processing. \citet{Kelly2005} develops Merleau-Ponty's view of the perceptual constancies, arguing that the real color or size of a thing is given as a norm towards which the varying appearances point, not as one appearance among them; we take this to describe how the series is organized by the Idea of the thing without being completed.

The third is the epistemology of perceptual fulfillment. \citet{Hopp2011} defends a phenomenological theory of perceptual knowledge on which perceptual experience is a fulfilling presentation of its object, and knowledge consists in the fulfillment of an empty intention by such a presentation. Hopp's account presupposes Principle~\ref{p:self}: if perception did not present the thing itself, the fulfillment of a thought by perception would not be the fulfillment of the thought by its object. It also shows why non-convergence is not a sceptical result: knowledge requires fulfillment of determinate intentions, not complete givenness. \citet{Drummond1990} had earlier argued that the Husserlian noema is not an intermediary between the act and the thing but the thing itself as intended, which is the same thesis in the vocabulary of the noema. \citet{Zahavi2003} presents the theory of adumbrations as the core of Husserl's perceptual realism and stresses that the inadequacy of thing-perception is an eidetic, not an empirical, feature of it.

\section{Objections}\label{s:objections}

\begin{objection}[The Kantian Idea]
Husserl holds that complete givenness of the thing is prescribed as an Idea in the Kantian sense, and that the infinite continuum of appearances is determinate as a whole. A determinate infinite whole of which every finite course is an initial segment is a convergent series.
\end{objection}

\begin{reply}
What Husserl calls determinate is the rule of the continuum, not a limit of givenness. The Idea is given with evidence as the Idea of an infinity of appearances, and its being given in this way is compatible with, indeed requires, that no course of appearances gives the thing completely \citep[§143]{Husserl1913}. Convergence in the sense of Definition~\ref{d:series} concerns what is given in fullness at each phase, and the Idea is not given in fullness at any phase. It is given as an Idea: the thing's complete givenness is itself given, but only as a task. The completion reading of the doctrine requires that the series accumulate, and Section~\ref{s:argument} shows that it does not.
\end{reply}

\begin{objection}[Science]
Physics determines a thing completely by specifying its physical state. Perception may not converge, but science reaches the complete thing.
\end{objection}

\begin{reply}
A theoretical determination is not a givenness of the thing. Physics determines a thing by idealized magnitudes whose sense is founded on perceptual measurement, and the founded character of science does not make its determinations intuitively given (PHE-17). Moreover, physical theory leaves open infinitely many determinations of the thing as it is perceived: its look under each illumination, its feel to each hand. The completeness of a physical state description is completeness with respect to the theory, not with respect to the thing as given.
\end{reply}

\begin{objection}[The sum of the series]
Sartre holds that the object is the infinite series of its appearances. If so, the thing is given as soon as the series is given, and the series can be given as a whole by a concept, as the series of natural numbers is.
\end{objection}

\begin{reply}
A concept of the series is not a givenness of its members. The series of natural numbers is given as a whole in the sense that its rule is grasped; its members are not thereby intuited. In the same way the rule of the series of appearances is grasped in each perception, which is why the thing is given as one and the same through its profiles. Grasping the rule is not giving the determinations in fullness.
\end{reply}

\begin{objection}[Scepticism]
If the thing is never given completely, we never perceive the thing itself, only its profiles; the withdrawal view is correct.
\end{objection}

\begin{reply}
The inference moves from the incompleteness of givenness to the absence of the thing, and Lemma~\ref{l:return} blocks it. What is given in a profile is the thing from a side, not the side as a separate object. Incomplete givenness of the thing is still givenness of the thing, just as non-primordial givenness of another's experience is still givenness of that experience (PHE-12). The withdrawal view needs a premise that our framework rejects: that a thing given incompletely is not given.
\end{reply}

\begin{objection}[Finite determinations]
Principle~\ref{p:open} begs the question. If the thing had finitely many determinations, a sufficiently large capacity would allow convergence.
\end{objection}

\begin{reply}
The conditional is correct, and the melody shows that it has instances: an object with finitely many determinations, all within the retentional horizon, can be given completely. The question is whether spatial things are such objects. Every determination of a spatial thing is a determination at a scale and under conditions, and both scales and conditions form open ranges. The argument for Principle~\ref{p:open} is not that we cannot list the determinations, but that every determination given intuitively points, in its horizon, to further determinations of the same kind that the same thing could present. This is a structure of the givenness of spatial things, not a limitation of our list.
\end{reply}

\section{What the result does not establish}

The result concerns spatial things and objects whose determinations are infinitely many. It does not apply to lived experiences given in reflection, which do not adumbrate themselves, nor to ideal objects such as essences or numbers, whose givenness has a different structure. It does not show that perceptual knowledge is impossible or defective; it shows only that perceptual knowledge is not complete givenness. It does not determine the bound in Principle~\ref{p:finite} for any subject. And it does not settle whether something beyond the series of appearances, a transphenomenal being in Sartre's sense, grounds the series; the theorem is about givenness, and so about what is given to someone as something (PHE-02), not about what there is.

\section{Conclusion}

We can return to the things themselves. Each adumbration gives the thing itself, and every determination of the thing can be brought to givenness by a suitable course of exploration. But the return is never complete. A phase of perception holds a bounded number of determinations in fullness, a spatial thing has infinitely many, and what is given in fullness is not accumulated but held in a moving window. So no phase gives the thing completely, and the series of adumbrations does not converge on a complete givenness. Complete givenness is given only as an Idea: as the rule of a series that no perception completes.

\begin{thebibliography}{99}\small
\bibitem[Drummond(1990)]{Drummond1990} Drummond, J.~J. (1990). \emph{Husserlian Intentionality and Non-Foundational Realism: Noema and Object}. Kluwer.
\bibitem[Gurwitsch(1964)]{Gurwitsch1964} Gurwitsch, A. (1964). \emph{The Field of Consciousness}. Duquesne University Press.
\bibitem[Heidegger(1927)]{Heidegger1927} Heidegger, M. (1927). \emph{Sein und Zeit}. Niemeyer.
\bibitem[Hopp(2011)]{Hopp2011} Hopp, W. (2011). \emph{Perception and Knowledge: A Phenomenological Account}. Cambridge University Press.
\bibitem[Husserl(1901)]{Husserl1901} Husserl, E. (1900--1901). \emph{Logische Untersuchungen}. 2 vols. Niemeyer.
\bibitem[Husserl(1913)]{Husserl1913} Husserl, E. (1913). \emph{Ideen zu einer reinen Phänomenologie und phänomenologischen Philosophie. Erstes Buch}. Niemeyer.
\bibitem[Husserl(1952)]{Husserl1952} Husserl, E. (1952). \emph{Ideen zu einer reinen Phänomenologie und phänomenologischen Philosophie. Zweites Buch} (Husserliana IV, M. Biemel, Ed.). Nijhoff.
\bibitem[Husserl(1966)]{Husserl1966} Husserl, E. (1966). \emph{Zur Phänomenologie des inneren Zeitbewusstseins (1893--1917)} (Husserliana X, R. Boehm, Ed.). Nijhoff.
\bibitem[Husserl(1973)]{Husserl1973} Husserl, E. (1973). \emph{Ding und Raum: Vorlesungen 1907} (Husserliana XVI, U. Claesges, Ed.). Nijhoff.
\bibitem[Kelly(2005)]{Kelly2005} Kelly, S.~D. (2005). Seeing things in Merleau-Ponty. In T. Carman \& M.~B.~N. Hansen (Eds.), \emph{The Cambridge Companion to Merleau-Ponty}. Cambridge University Press.
\bibitem[Madary(2017)]{Madary2017} Madary, M. (2017). \emph{Visual Phenomenology}. MIT Press.
\bibitem[Merleau-Ponty(1945)]{MerleauPonty1945} Merleau-Ponty, M. (1945). \emph{Phénoménologie de la perception}. Gallimard.
\bibitem[Merleau-Ponty(1964)]{MerleauPonty1964} Merleau-Ponty, M. (1964). \emph{Le visible et l'invisible} (C. Lefort, Ed.). Gallimard.
\bibitem[Noë(2004)]{Noe2004} Noë, A. (2004). \emph{Action in Perception}. MIT Press.
\bibitem[Sartre(1943)]{Sartre1943} Sartre, J.-P. (1943). \emph{L'être et le néant: essai d'ontologie phénoménologique}. Gallimard.
\bibitem[Zahavi(2003)]{Zahavi2003} Zahavi, D. (2003). \emph{Husserl's Phenomenology}. Stanford University Press.
\end{thebibliography}

\end{document}
